\SourceLeaf{322}\GaussVSourceFolio{316}{L}{ERDMAGNETISMUS UND MAGNETOMETER.}
\noindent sie schlechthin fordern muss, aus ihnen aber die Erscheinungen erschöpfend vollständig als nothwendig ableitet.

Dagegen äussert die erdmagnetische Kraft uns. erkennbare Wirkungen nur auf einige Arten von Körpern, auf diejenigen nemlich, auf welche durch wirkliche Magnete, natürliche oder künstliche, gewirkt werden kann, also wenn wir die erst in der jüngsten Zeit entdeckte Wechselwirkung zwischen Magnetismus und galvanischen Strömen beiseite setzen, auf magnetische oder magnetisirbare Körper. Das weiche Eisen \textit{macht} die erdmagnetische Kraft magnetisch ohne Beharrlichkeit; hingegen einen schon mit beharrlichem Magnetismus versehenen Körper, sei es ein natürlicher Magnet, oder ein künstlicher aus gehärtetem Stahl, \textit{bewegt} die erdmagnetische Kraft nach bestimmten Gesetzen. Von der letztern Wirkung soll hier allein die Rede sein: die der Wirkung unterworfenen Träger eines beharrlichen Magnetismus, am besten von nadelförmiger oder länglich prismatischer Gestalt, sollen, welche Grösse sie auch haben mögen, Magnetnadeln heissen.

Durch die Richtung der Schwerkraft an jedem Orte wird die gerade Linie bestimmt, die wir eine Verticallinie nennen, und der Gegensatz des Oben und Unten. Die Astronomie lehrt uns, die Lage dieser Linie gegen den Erdäquator und gegen eine willkürlich gewählte Meridianebene bestimmen, und liefert dadurch die mathematischen Grundlagen der Geographie. Unsre feinsten Beobachtungen vermögen nicht, in der Richtung der Schwerkraft an einem gegebenen Orte auch nur die geringste Veränderung zu erkennen, obwohl wir aus theoretischen Gründen sehr wohl wissen, dass diese Richtung unaufhörlichen Veränderungen unterworfen sein muss. Denn die Schwerkraft ist nur die Gesammtwirkung aller Theile des Erdkörpers, etwas modificirt durch die Centrifugalkraft vermöge der Rotationsbewegung, und durch die fremden Weltkörper: allein die ganze letztere unmittelbare Wirkung auf die Schwerkraft, und die mittelbare, durch die beständigen Bewegungen vieler Kubikmeilen von Wassermassen vermöge der Ebbe und Fluth, bleibt viel zu klein für das Messungsvermögen unsrer feinsten Instrumente; noch mehr verschwindet also die Wirkung von sonstigen Versetzungen von Massen auf der Erdoberfläche durch andere Natur- oder Menschenkräfte.

Ganz anders verhält es sich in dieser Beziehung mit der Richtung der erdmagnetischen Kraft. Scharf in sich bestimmt ist auch sie an jedem Orte, aber,
\par\clearpage
\SourceLeaf{323}\GaussVSourceFolio{317}{R}{ERDMAGNETISMUS UND MAGNETOMETER.}
\noindent genau zu reden, nur in jedem Augenblick. Wir beziehen diese Richtung an jedem Orte auf die Verticallinie (oder, was auf dasselbe hinausläuft, auf die gegen diese normale Horizontalebene) und auf die Meridianebene. Den Winkel, welchen die Richtung der erdmagnetischen Kraft mit der Horizontalebene macht, nennen wir die Neigung (Inclination) der Magnetnadel; der Winkel zwischen derjenigen Verticalebene, in welcher sich jene Richtung befindet, und der Meridianebene ist die Abweichung (Declination) der Nadel. Diese beiden Elemente bestimmen die Richtung der erdmagnetischen Kraft vollständig: sie sind an verschiedenen Orten verschieden, aber sie sind an einem und demselben Orte nicht beständig, sondern immerwährenden Veränderungen unterworfen, auf die wir nachher zurückkommen werden.

Zunächst müssen wir aber die ungleiche Art, \textit{wie} die beiden Kräfte nach ihren Richtungen wirken, näher betrachten. Die Schwerkraft treibt einen Körper, sobald keine Hindernisse im Wege stehen, in ihrer Richtung nach unten, und diese Bewegung wird immer schneller, so lange der Körper frei fallen kann. Die Schwerkraft bringt einen Körper, der sich frei bewegen kann, nur in eine fortschreitende (progressive) Bewegung, nicht in eine drehende (rotatorische).

Mit der erdmagnetischen Kraft verhält es sich gerade umgekehrt: diese kann den Körpern, welche sie in Bewegung setzt, nur eine drehende, nie eine fortschreitende Bewegung ertheilen. Wollen wir also die Wirkung der erdmagnetischen Kraft auf einen Körper rein beobachten, so müssen wir zuvörderst die progressiven Bewegungen, die die Schwere hervorbringen könnte, ausschliessen oder unmöglich machen. Eine Magnetnadel, die auf ihrer untern Seite eine kegelförmige Vertiefung (ein Hütchen) hat, und damit auf einer feinen Spitze hängt, befindet sich in diesem Falle. Trifft dieser Aufhängepunkt (die Spitze des Hütchens) genau mit dem Schwerpunkt der Nadel zusammen, so ist letztere gegen die Schwerkraft ganz indifferent, und zeigt sich nun einer Kraft unterwürfig, die sie in die Richtung des Erdmagnetismus zu bringen strebt. Ist die Nadel schon Anfangs in dieser Lage, so bleibt sie darin: trifft aber der Erdmagnetismus sie Anfangs in einer andern Lage, so setzt letzterer sie in eine Bewegung, vermöge welcher sie sich jener Lage nähert, und (weil die Schnelligkeit der Bewegung so lange zunimmt, als jene Lage noch nicht erreicht ist) sogar über dieselbe hinausgeht, wo dann aber die erdmagnetische Kraft, stets die Nadel der Normalrich-
\par\clearpage
\SourceLeaf{324}\GaussVSourceFolio{318}{L}{ERDMAGNETISMUS UND MAGNETOMETER.}
\noindent tung näher zu bringen strebend, die Schnelligkeit der Bewegung fortwährend wieder vermindert, bis diese vernichtet ist, und rückwärts geht. Auf diese Weise macht die Nadel \textit{Schwingungen}, desto grössere, je mehr die ursprüngliche Lage von der Normalrichtung abwich, und die Normalrichtung liegt in oder nahe bei der Mitte der Schwingungen. Das erstere würde genau der Fall sein, wenn nicht äussere Hindernisse, der Widerstand der Luft und die Reibung im Hütchen die Bewegung nach und nach lähmten: diese Hindernisse vermindern nach und nach die Grösse des Schwingungsbogens, bis zuletzt die Nadel in der Richtung der erdmagnetischen Kraft zur Ruhe kommt.

Man pflegt jedoch die Spitze des Hütchens, oder den Aufhängepunkt, um welchen die Nadel sich frei bewegen kann, nicht im Schwerpunkt der Nadel, sondern etwas höher, anzubringen, wodurch sich die Erscheinung etwas anders gestaltet. Es entsteht dann ein Conflict der Schwerkraft mit der erdmagnetischen Kraft, und die Nadel stellt sich nicht mehr genau in die Richtung der letztern Kraft, aber ihr so nahe, wie es dieser Conflict verstattet. Die Schwerkraft strebt nemlich, den Schwerpunkt senkrecht unter den Aufhängungspunkt zu bringen; bei der Stellung der Nadel, genau nach der Richtung der erdmagnetischen Kraft, würde aber der Schwerpunkt einen etwas höhern Platz erhalten (wenn nicht zufällig die gegenseitige Lage beider Punkte in der Nadel schon die zu jener Richtung erforderliche ist): die Natur, stets die distributive Gerechtigkeit auf das strengste verwaltend, ertheilt daher der Nadel, nach Maassgabe der Stärke beider Kräfte, eine vermittelnde Zwischenlage, wobei von der genauen Inclination weniger oder mehr aufgeopfert werden muss, die aber nothwendig mit dem magnetischen Meridian, d. i. derjenigen Verticalebene, in der sich die eigentliche Richtung der erdmagnetischen Kraft befindet, übereinstimmt. Wie allen Geschäften eine verständige Theilung stets zum Vortheil gereicht, so trennt auch der Naturforscher die Ausmittelung der Declination von der Inclination, und hängt, wo es ihm zunächst um erstere zu thun ist, seine Nadel nicht im Schwerpunkt auf, sondern so, dass eben die Declination am besten hervortritt: er hängt sie so auf, dass sie horizontal schwebt. Der Seefahrer erreicht dieses, indem er, wenn er in den Bereich einer beträchtlich geänderten Inclination kommt, seine Nadel auf der einen Seite mit einem leichten, gegen den Magnetismus indifferenten Körperchen, z. B. einem Stückchen Wachs, belastet. Der Naturforscher, der für feinere Zwecke die Nadel gar nicht mit einem Hütchen, sondern an ei-
\par\clearpage
\SourceLeaf{325}\GaussVSourceFolio{319}{R}{ERDMAGNETISMUS UND MAGNETOMETER.}
\noindent nem feinen Faden aufhängt, legt sie in ein an das untere Ende des Fadens geknüpftes Schiffchen, in welchem er sie, so viel zu obigem Zweck nöthig, verschiebt, oder auch sie mit einem leichten Laufgewicht belastet.

Die Wirkungsart unsrer beiden Naturkräfte stellt sich hienach wesentlich verschieden dar: es wird aber interessant sein, zu zeigen, wie die Schwerkraft unter geeigneten Umständen ganz analoge Wirkungen hervorbringen kann. Die Hydrostatik lehrt, dass ein in eine Flüssigkeit eingetauchter Körper so viel an seinem Gewicht verliert, als das Gewicht der verdrängten Flüssigkeit austrägt. Ein Körper, specifisch schwerer als Wasser, wird im Innern von Wasser auch noch in der Verticallinie nach unten getrieben, aber nur mit einer verhältnissmässig geringern Kraft; ein specifisch leichterer Körper hingegen nach oben, und steht also gleichsam unter dem Einfluss einer negativen Schwere; endlich ein fester Körper, genau von demselben specifischen Gewicht wie Wasser, wird weder nach unten, noch nach oben getrieben, sondern erhält sich in der Höhe, in welcher er in Ruhe sich einmal befindet. Sind diese Körper homogen, so erhalten sie in den beiden ersten Fällen auch nur progressive Bewegungen (insofern wir von dem Widerstande abstrahiren), und in dem letzten Fall verhält sich der Körper inmitten des Wassers völlig indifferent, als ob die Schwere für ihn gar nicht da wäre. Anders verhält es sich aber mit einem Körper, der aus Theilen von ungleichem specifischem Gewicht zusammengesetzt ist. Denken wir uns einen länglichen prismatischen Stab, dessen eine Hälfte von Elfenbein, die andere von Kork ist. Das specifische Gewicht des Elfenbeins übertrifft das des Wassers wenig mehr, als das specifische Gewicht des Korks gegen letzteres zurücksteht. Wir setzen diese kleine Ungleichheit des Unterschiedes hier bei Seite, oder denken uns, anstatt reinen Wassers, mit Salz soweit versetztes, dass das specifische Gewicht der Flüssigkeit genau mitten inne steht zwischen den specifischen Gewichten der beiden festen Theile des Stabs. Dieser Stab, im Innern einer Wassermasse, wird nun, da sein specifisches Gewicht im Ganzen dasselbe ist, wie das des Wassers, weder nach unten, noch nach oben, aber das Elfenbeinende wird nach unten, das Korkende nach oben getrieben: der Stab erhält keine progressive, wohl aber eine drehende Bewegung, wenn er sich nicht schon Anfangs in senkrechter Lage befand. Dasselbe wird auch noch gelten, wenn dieser Stab mit einer Holzart überlegt ist, die das specifische Gewicht des Wassers hat, oder auch, wenn in einem Stab aus solchem Holz an dem einen Ende ein Stück El-
\par\clearpage
\SourceLeaf{326}\GaussVSourceFolio{320}{L}{ERDMAGNETISMUS UND MAGNETOMETER.}
\noindent fenbein, an dem andern ein eben so grosses Stück Kork eingelegt ist. Wir haben demnach hier ganz das Gegenstück von der Wirkungsart der erdmagnetischen Kraft, nur dass an die Stelle der dieser eigenthümlichen Richtung jetzt die Verticallinie getreten ist, und werden dadurch auf eine Vorstellungsart geführt, die zur Erklärung der Wirkung der erdmagnetischen Kraft auf die Magnetnadel dient. Wir nehmen nemlich zwei Stoffe an, auf welche diese Kraft unmittelbar auf ähnliche Art wirkt, wie die Schwerkraft auf alle ponderabeln Körper, indem sie dieselben, zwar in Einer bestimmten geraden Linie, aber in entgegengesetzten Richtungen, zu bewegen strebt. Diese beiden Stoffe müssen wir an die Magnetnadel fest gebunden voraussetzen (weil sonst die erdmagnetische Kraft nur die Stoffe \textit{in} der Nadel, nicht diese selbst, bewegen würde); den einen an das eine Ende, den andern an das andere, und die Quantität des einen Stoffs muss in jeder Nadel der Quantität des andern genau gleich sein (weil sonst auch eine progressive Bewegung erfolgen müsste). Man nennt diese Stoffe magnetische \textit{Fluida}, um ihre leichte Beweglichkeit in dem sich nicht zu beharrlicher Magnetisirung eignenden weichen Eisen zu bezeichnen, und unterscheidet sie durch die Zusätze nördliches und südliches Fluidum, indem dasjenige Ende der Nadel, welches das erstere trägt (der Nordpol), sich an den meisten Orten der Erde nach der Nordseite richtet. Das nördliche Fluidum pflegt man auch das positive, das südliche das negative zu nennen. Der Stahl ist dabei nur der Träger dieser Fluida, wie in dem vorhergehenden Gleichniss die Holzumgebung der Träger des Elfenbeins und Korks, und die erdmagnetische Kraft wirkt auf jenen nur mittelbar.

Diese Vorstellungsart bedarf aber noch einer Modification. Unser Elfenbein-Kork-Stab würde offenbar noch dieselbe Erscheinung (wenn auch in geringerer Stärke) darbieten, wenn er, anstatt Ein Stück Elfenbein und Ein eben so grosses Stück Kork zu enthalten, aus mehrern, immerhin auch sehr vielen, Paaren zusammengesetzt wäre, vorausgesetzt, dass diese Paare in gehöriger Ordnung liegen. Bei unserer ersten Voraussetzung würde der Stab seine Eigenschaft verlieren, wenn man ihn in der Mitte zerschnitte; bei der andern hingegen bleibt die Eigenschaft nach jeder Zerschneidung, wo nur kein zusammengehöriges Paar getrennt wird.

Die Erfahrung ergibt, dass, wenn ein Magnetstab in der Mitte, oder an irgend einer andern Stelle durchgebrochen wird, beide Stücke sich sogleich wieder als Magnete zeigen, die von der erdmagnetischen Kraft nur eine drehende, nie
\par\clearpage
\SourceLeaf{327}\GaussVSourceFolio{321}{R}{ERDMAGNETISMUS UND MAGNETOMETER.}
\noindent eine progressive Bewegung erhalten. Wir sind daher genöthigt, anzunehmen, dass in der Magnetnadel die magnetischen Fluida zwar getrennt sind, aber nicht so, dass das eine Fluidum sich am einen, das andere am andern Ende befinde, sondern vielmehr so, dass wir die Nadel wie ein Aggregat von unzähligen für uns unmessbar kleinen Stahltheilchen betrachten müssen, deren jedes eben so viel nördliches wie südliches Fluidum im getrennten Zustande enthält.

Wir haben bisher den Magnet nur in Beziehung auf diejenige Wirkung betrachtet, welche die erdmagnetische Kraft auf ihn ausübt, weil diese zunächst den Gegenstand des gegenwärtigen Aufsatzes ausmacht: viel länger war schon diejenige Wirkung bekannt, welche zwei Magnete auf einander ausüben, und die in einer gegenseitigen Anziehung der ungleichnamigen Pole und einer Abstossung der gleichnamigen besteht. Nach Beschaffenheit der Umstände können dadurch drehende und fortschreitende Bewegungen erregt werden. Es bedarf zur Erklärung dieser Phänomene nichts weiter, als noch anzunehmen, dass die magnetischen Fluida auf einander wirken, die gleichnamigen abstossend, die ungleichnamigen anziehend, und wir wissen jetzt aus scharfen Versuchen, dass die Stärke dieser Abstossung oder Anziehung zwischen zweien Theilchen solcher Flüssigkeiten eben so im umgekehrten Verhältniss des Quadrats der Entfernung steht, wie die allgemeine gegenseitige Anziehung aller ponderabeln Körper.

Nur kurz erwähnen wir endlich (da es nicht unmittelbar zu unserm gegenwärtigen Zweck gehört) der Wirkung der Magnete auf nicht magnetisirten Stahl und weiches Eisen, welche Wirkung bekanntlich in einer Anziehung besteht. Sie ist eine Folge des eben angeführten Verhaltens der magnetischen Flüssigkeiten, die in allem Stahl und Eisen als schon vorhanden angesehen werden müssen, und bei der Annäherung eines Magnets eine Scheidung erleiden, so dass jene Körper dadurch selbst magnetisch werden. Nur ist das weiche Eisen für sich nicht fähig, die Trennung der magnetischen Fluida in seinem Innern dauernd zu erhalten. Ein Stück weiches Eisen, mit einem Ende an einem Magnet hängend (oder ihm auch nur nahe gebracht), verhält sich so lange selbst wie ein Magnet, verliert aber nach der Trennung oder Entfernung von jenem diese Eigenschaft nach wenigen Augenblicken fast ganz wieder, während ein Stück gehärteten, aber noch nicht magnetisirten Stahls (in welchem die Trennung der magnetischen Flüssigkeiten schwerer geschieht, aber, einmal erfolgt, viel bleibender ist), theils überhaupt von einem Magnet schwächer angezogen wird, als weiches Eisen, theils
\par\GaussVSourceGathering{41}
\par\clearpage
\SourceLeaf{328}\GaussVSourceFolio{322}{L}{ERDMAGNETISMUS UND MAGNETOMETER.}
\noindent auch nach der Trennung den Grad von Magnetismus, welchen es in jener Verbindung erhalten hat, auf längere Zeit beibehält.

Wir kehren zu der erdmagnetischen Kraft zurück, deren Kenntniss erst vollständig wird, wenn man ausser ihrer Richtung auch ihre Stärke (Intensität) angeben kann. Um diese auszumessen, ist man auf folgende Art zu Werke gegangen.

Wir haben oben gesehen, unter welchen Umständen eine aufgehängte Magnetnadel in Schwingungen versetzt wird: erwägen wir jetzt näher, von welchen Umständen \textit{die Dauer} einer Schwingung abhängt.

Zuerst ist diese Dauer, alles übrige gleich gesetzt, von der Grösse des Schwingungsbogens abhängig; jene ist desto kleiner, je kleiner dieser ist, so jedoch, dass bei immer abnehmenden Bögen die Dauer sich nur immer mehr einem Grenzwerthe nähert, ohne in mathematischer Schärfe solchen erreichen zu können. Das Verhältniss der Schwingungsdauer für jede Grösse des Schwingungsbogens zu dem Grenzwerthe kann man durch bekannte Formeln berechnen. Für eine Nadel z. B., welche einen Schwingungsbogen von 180 Grad in 23.6068 Secunden beschreibt, ist der Grenzwerth der Schwingungsdauer 20 Secunden, und folgende Übersicht gibt eine Vorstellung von der successiven Annäherung zu demselben für immer kleinere Bögen:

\begin{center}
\newcommand{\GaussVSourceOscillationRow}[4]{%
  \makebox[2em][r]{#1}\hspace{.4em}\makebox[2.9em][c]{#2}&%
  \makebox[3.6em][l]{#3}\hspace{.4em}\makebox[4.7em][c]{#4}%
}
\begin{tabular}{@{}r@{\hspace{1.2em}}|@{\hspace{1.1em}}l@{}}
\textit{Schwingungsbogen.}&\textit{Schwingungsdauer.}\\
\GaussVSourceOscillationRow{180}{Grad.}{23.6068}{Secunden.}\\
\GaussVSourceOscillationRow{120}{„}{21.4636}{„}\\
\GaussVSourceOscillationRow{60}{„}{20.3482}{„}\\
\GaussVSourceOscillationRow{30}{„}{20.0860}{„}\\
\GaussVSourceOscillationRow{20}{„}{20.0381}{„}\\
\GaussVSourceOscillationRow{10}{„}{20.0095}{„}\\
\GaussVSourceOscillationRow{8}{„}{20.0061}{„}\\
\GaussVSourceOscillationRow{6}{„}{20.0034}{„}\\
\GaussVSourceOscillationRow{4}{„}{20.0015}{„}\\
\GaussVSourceOscillationRow{2}{„}{20.0004}{„}\\
\GaussVSourceOscillationRow{1}{„}{20.0001}{„}\\
\end{tabular}
\end{center}

Man sieht daraus, dass bei Schwingungsbögen von mässiger Grösse der Unterschied einer Schwingungsdauer von dem Grenzwerthe kaum merklich ist.
\par\clearpage
\SourceLeaf{329}\GaussVSourceFolio{323}{R}{ERDMAGNETISMUS UND MAGNETOMETER.}
\noindent Dieser Grenzwerth wird immer verstanden, so oft man von Schwingungsdauer schlechthin spricht, und in der herkömmlichen mathematischen Sprache als Schwingungsdauer für einen unendlich kleinen Bogen bezeichnet.

Zweitens hängt die Schwingungsdauer einer Nadel ab von der Stärke ihrer Magnetisirung. Behandelt man eine Anfangs schwach magnetisirt gewesene Nadel mit kräftigern Streichmitteln, so werden ihre Schwingungen schneller. Es gibt jedoch für jede Nadel einen bestimmten höchsten Grad von Magnetismus, den sie annehmen oder festhalten kann, und den man wohl die Sättigung nennt. Allein es ist von Wichtigkeit, hier zu bemerken, dass die Bestimmung dieses Sättigungspunkts der Erfahrung zufolge einer sehr grossen Schärfe nicht fähig zu sein scheint. Wenn man eine und dieselbe Nadel in öftern Wiederholungen auch mit den kräftigsten Mitteln magnetisirt, nachdem man dazwischen ihr den Magnetismus zum Theil wieder entzogen hatte, so geben doch die jedesmaligen Schwingungszeiten keineswegs einen \textit{solchen} Grad von Übereinstimmung, als man für Normalbestimmungen fordern müsste.

Der dritte Umstand, welcher die Schwingungsdauer bestimmt, ist die Grösse der Nadel. Von zwei ungleich grossen Nadeln, die jede in ihrer Art gleich gut magnetisirt sind, wird die grössere langsamer schwingen. Grössere Dicke und Breite, so lange diese Dimensionen gegen die Länge noch sehr klein bleiben, hat dabei einen geringern Einfluss, als vergrösserte Länge. Eine grosse gut magnetisirte Nadel hat zwar einen stärkern Magnetismus, als eine kleinere, ja wenn beide ähnliche Gestalten haben, so wird man sie nur dann gleich gut in ihrer Art magnetisirt nennen können. wenn das Verhältniss des Magnetismus dasselbe ist, wie das der Grösse (dem Raume oder Gewichte nach): dass dann, dessen ungeachtet, die grössere langsamer schwingt, ist eine nothwendige Folge davon, dass der stärkere Magnetismus nicht bloss grössere Masse zu bewegen, sondern durch grössere Räume zu bewegen hat.

Viertens endlich ist die Schwingungsdauer abhängig von der Stärke der erdmagnetischen Kraft selbst. In der That muss eine gegebene Nadel, in bestimmter Stärke magnetisirt, schnellere oder langsamere Schwingungen machen, je nachdem die auf sie wirkende erdmagnetische Kraft stärker oder schwächer ist, und es bietet sich also ein Mittel dar, die Stärke dieser Kraft an verschiedenen Orten zu vergleichen, indem man \textit{eine und dieselbe} Nadel daselbst schwingen lässt. Man weiss, dass eine doppelt schnellere Schwingung einer vierfach grössern Kraft,
\par\GaussVSourceGathering{41\textsuperscript{*}}
\par\clearpage
\SourceLeaf{330}\GaussVSourceFolio{324}{L}{ERDMAGNETISMUS UND MAGNETOMETER.}
\noindent eine dreifach schnellere Schwingung einer neunfach grössern Kraft entspricht u. s. w., so dass die Quadratzahl von der Anzahl der Schwingungen in einer beliebig gewählten Zeit, z. B. in einer Minute, als das Maass der Kraft angesehen werden kann. Übrigens ist hier immer die erdmagnetische Kraft zu verstehen, so weit sie in derjenigen Ebene wirkt, in welcher die Schwingung geschieht, mithin die ganze erdmagnetische Kraft, wenn die in ihrem Schwerpunkt aufgehängte Nadel in der Ebene des magnetischen Meridians schwingt, hingegen nur der horizontale Theil der erdmagnetischen Kraft, wenn die Nadel, oberhalb ihres Schwerpunkts aufgehängt, Schwingungen in horizontaler Ebene macht. Schwingungen der letztern Art lassen sich viel bequemer und schärfer beobachten, als die der erstern, und für die Anwendung sind jene eben so brauchbar, da das Verhältniss der ganzen erdmagnetischen Kraft zu ihrem horizontalen Theile auf bekannte Weise von der Inclination abhängt.

Man hat auf diese Weise die Intensitäten der erdmagnetischen Kraft an vielen Örtern der Erde unter einander verglichen, indem man auf Reisen, die zum Theil hauptsächlich zu diesem Zweck unternommen waren, eine oder mehrere Nadeln mit sich führte, und deren Schwingungsdauer beobachtete: als Einheit für die Resultate kann man die Stärke der erdmagnetischen Kraft an einem beliebig gewählten Orte annehmen.

Offenbar ist dies Verfahren ganz von der Voraussetzung abhängig, dass der magnetische Zustand der angewandten Nadeln ganz ungeändert bleibt. Allein es gibt mehrere Ursachen, die diesen Zustand verändern können. Zuvörderst hat die Temperatur einen sehr merklichen Einfluss auf diesen Zustand. Bei grösserer Wärme wird der Magnetismus einer Nadel schwächer, kommt jedoch mit dem frühern Temperaturzustande wieder auf seine vorige Stärke zurück, wenn die Wärme innerhalb mässiger Grenzen geblieben ist. \textit{Diese} Veränderlichkeit kann man daher durch Rechnung berücksichtigen und unschädlich machen; vor zu starker Erhitzung muss man aber die Nadel wohl in Acht nehmen, weil dadurch ihr Magnetismus bleibend geschwächt wird. Ferner darf man die Nadel nicht mit andern Magneten oder auch mit Eisen in Berührung bringen, weil man sonst nach der Abtrennung durchaus nicht darauf rechnen kann, den vorigen magnetischen Zustand der Nadel genau wieder zu finden. Allein auch bei aller solcher Vorsicht hat man doch für völlige Unwandelbarkeit dieses Zustandes keine Bürgschaft. Nadeln aus schwach gehärtetem Stahl verlieren schon in kurzer Zeit ei-
\par\clearpage
\SourceLeaf{331}\GaussVSourceFolio{325}{R}{ERDMAGNETISMUS UND MAGNETOMETER.}
\noindent nen beträchtlichen Theil ihres Magnetismus; gut gehärtete halten ihn besser an sich; allein auch bei den bestgehärteten wird man immer im Laufe der Zeit einige Veränderung zu befürchten haben. Man könnte glauben, dass diese Veränderung sich an der veränderten Schwingungsdauer der Nadel an einem und demselben Orte erkennen lasse, und in Beziehung auf beträchtliche und schon nach mässigen Zeitintervallen eingetretene Veränderungen ist dies auch ganz richtig: allein dieser Schluss würde ganz illusorisch sein, wenn man ihn auf kleine Veränderungen oder auf sehr lange Zeiträume ausdehnen wollte: denn so wie der Erfahrung zufolge die Richtung der erdmagnetischen Kraft an einem Orte sehr grossen Veränderungen unterworfen ist, wird dies ohne Zweifel auch mit der Intensität dieser Kraft der Fall sein, daher man, wenn man nach einer Reihe von Jahren die Nadel an einem Orte andere Schwingungen machen sieht, als früher, völlig im Dunkeln bleibt, wieviel Antheil daran die Veränderung der Nadel, und wieviel die Veränderung der Stärke des Erdmagnetismus gehabt hat. Das Resultat dieser Betrachtungen ist also, dass die erwähnte comparative Methode sehr nützliche Dienste leistet, wenn man sie nur auf Bestimmungen innerhalb eines mässigen Zeitraumes anwendet, und es an der nöthigen Vorsicht nicht fehlen lässt; dass jedoch die Zuverlässigkeit und Genauigkeit dieses Verfahrens immer nur beschränkt bleibt, und dass die so hoch interessante Frage, welchen Veränderungen die Intensität der erdmagnetischen Kraft an einem Orte im Laufe langer Zeiträume unterworfen sein mag, auf diese Weise gar nicht zu beantworten ist.

Wir machen uns frei von jener Unsicherheit, und gewinnen das Mittel zur Beantwortung dieser Frage, indem wir an die Stelle der comparativen Methode eine andere setzen, die die Intensität des Erdmagnetismus auf ein von der Individualität der gebrauchten Magnetnadeln ganz unabhängiges absolutes Maass zurückführt, d. i. auf ein solches, welches nur auf den für sich feststehenden, jederzeit mit äusserster Schärfe wieder nachzuweisenden Raum-, Zeit- und Gewichtseinheiten beruhet.

Zu einer vollkommenen Einsicht in das Wesen dieser Methode würde eine viel ausführlichere Entwicklung nothwendig sein, als hier Platz finden kann, zumal unter Verzichtleistung auf eine mit Wenigem viel sagende mathematische Einkleidung. Indessen wird die folgende Darstellung wenigstens die Möglichkeit der Zurückführung der Stärke der erdmagnetischen Kraft auf absolutes Maass begreiflich machen. Da es nach der schon oben gemachten Bemerkung nur auf
